\chapter{Analiza porównawcza}
\label{ch:analiza}

Niniejszy rozdział zawiera statystyczną i~jakościową analizę wyników eksperymentów
opisanych w~rozdziale~\ref{ch:wyniki}. Miarą jakości frontów Pareto jest
hiperobjętość (ang. \textit{hypervolume}, HV), a~wszystkie wnioski oparto na
testach statystycznych dla prób sparowanych po instancjach problemu.
W~podrozdziale \ref{sec:analiza_metodyka} opisano metrykę, sposób jej normalizacji,
protokół eksperymentalny oraz zastosowane testy. Podrozdział
\ref{sec:analiza_hv} przedstawia główne wyniki porównania hiperobjętością,
a~podrozdział \ref{sec:analiza_pary} --- testy dla par algorytmów wraz z~wielkością
efektu. W~podrozdziale \ref{sec:analiza_kryteria} przeanalizowano zachowanie
algorytmów w~przestrzeni poszczególnych kryteriów, w~podrozdziale
\ref{sec:analiza_zbieznosc} --- przebieg zbieżności, a~w~podrozdziale
\ref{sec:analiza_budzet} --- wrażliwość na budżet obliczeniowy i~czas obliczeń.
Podrozdział \ref{sec:analiza_interpretacja} zawiera interpretację zaobserwowanych
różnic w~odniesieniu do mechanizmów algorytmów, podrozdział
\ref{sec:analiza_niestacjonarnosc} --- ocenę przydatności metod w~warunkach
niestacjonarnych. Rozdział zamyka omówienie ograniczeń przeprowadzonej analizy
(podrozdział~\ref{sec:analiza_ograniczenia}); zbiorcze wnioski z~pracy
przedstawiono w~rozdziale~\ref{ch:podsumowanie}.

\section{Metodyka porównania}
\label{sec:analiza_metodyka}

\subsection{Hiperobjętość i~jej normalizacja}
\label{sec:analiza_hv_def}
Hiperobjętość zbioru rozwiązań $F$ względem punktu referencyjnego $\mathbf{r}$
to miara Lebesgue'a obszaru przestrzeni kryteriów zdominowanego przez $F$
i~ograniczonego przez $\mathbf{r}$:
\begin{equation}
    \label{eq:hv}
    \mathrm{HV}(F, \mathbf{r}) = \Lambda\Bigl(\,\bigcup_{\mathbf{f} \in F}
        \bigl[f_1, r_1\bigr] \times \bigl[f_2, r_2\bigr] \times \bigl[f_3, r_3\bigr]\Bigr)
\end{equation}
Metryka ta jest zgodna z~relacją dominacji Pareto (front dominujący inny front ma
nie mniejszą hiperobjętość) i~jako jedyna z~powszechnie stosowanych miar jakości
jednocześnie nagradza zbieżność do frontu, jego rozpiętość i~równomierność
rozkładu rozwiązań, nie wymagając znajomości frontu optymalnego --- co w~rozważanym
problemie jest warunkiem koniecznym, gdyż front optymalny nie jest znany.

Kryteria problemu mają nieporównywalne skale (makespan rzędu setek tików, energia
rzędu tysiąca jednostek, ryzyko radarowe od zera do kilkudziesięciu tysięcy), więc
surowa hiperobjętość byłaby zdominowana przez kryterium o~największym zakresie.
Dla każdej instancji problemu wartości kryteriów znormalizowano zatem do wspólnych
granic wyznaczonych przez sumę frontów końcowych wszystkich trzech algorytmów:
\begin{equation}
    \label{eq:hv_norm}
    f_m^{n} = \frac{f_m - z_m^{\min}}{z_m^{\max} - z_m^{\min}},
    \qquad
    z_m^{\min} = \min_{\mathbf{f} \in F_{\cup}} f_m,\;\;
    z_m^{\max} = \max_{\mathbf{f} \in F_{\cup}} f_m,\;\;
    F_{\cup} = F_{\text{NSGA-II}} \cup F_{\text{NSGA-III}} \cup F_{\text{MOGWO}}
\end{equation}
Punkt referencyjny przyjęto jako $\mathbf{r} = (1{,}1;\, 1{,}1;\, 1{,}1)$, tj.
z~marginesem $10\%$ poza najgorszą wartością każdego kryterium, dzięki czemu
skrajne rozwiązania frontu wnoszą niezerowy wkład. Wyniki podawane są jako iloraz
$\mathrm{HV}/\mathrm{HV}_{\max}$, gdzie $\mathrm{HV}_{\max} = 1{,}1^3 = 1{,}331$
jest objętością całego prostopadłościanu referencyjnego; wartość $1$ oznaczałaby
front dominujący cały obszar rozpięty przez sumę frontów. Tak znormalizowana
hiperobjętość jest miarą \emph{względną}: porównuje algorytmy w~obrębie jednej
instancji, natomiast jej wartości w~różnych instancjach nie są bezpośrednio
porównywalne. Z~tego powodu wszystkie testy statystyczne przeprowadzono dla prób
sparowanych po instancjach.

Hiperobjętość obliczano dokładnie (bez próbkowania Monte Carlo) algorytmem
rekurencyjnego cięcia wymiarów. Poprawność implementacji zweryfikowano testami
jednostkowymi względem wartości wyznaczonych ze wzoru włączeń--wyłączeń oraz
niezależną implementacją w~języku Python użytą w~skrypcie analizy --- obie
implementacje dały wartości zgodne z~dokładnością $4 \cdot 10^{-7}$ na wszystkich
porównanych frontach.

\subsection{Protokół eksperymentalny}
\label{sec:analiza_protokol}
Eksperyment główny objął $20$ niezależnych instancji problemu (teren generowany
z~ziarnami $42$--$61$, scenariusz domyślny z~tabeli~\ref{tab:scenariusz_parametry}).
Na każdej instancji uruchomiono wszystkie trzy algorytmy z~identycznym budżetem:
populacja (wataha) $N = 100$ i~$500$ iteracji, czyli $50\,000$ ocen funkcji celu na
algorytm. Ziarno generatora liczb pseudolosowych algorytmu wyznaczano
deterministycznie z~ziarna instancji, co zapewnia pełną odtwarzalność. Dodatkowo,
w~celu zbadania wrażliwości na budżet, na pierwszych $10$ instancjach powtórzono
porównanie dla $100$, $250$ i~$1000$ iteracji (wyniki dla $500$ iteracji pochodzą
z~pierwszych $10$ instancji eksperymentu głównego). W~trakcie każdego uruchomienia
co $10$ iteracji zapisywano migawkę bieżącego zbioru niezdominowanego (pierwszy
front populacji w~NSGA-II/III, archiwum w~MOGWO), oceniając ją później w~tych samych
granicach normalizacji co front końcowy --- z~tych migawek wyznaczono krzywe
zbieżności.

\subsection{Testy statystyczne}
\label{sec:analiza_testy}
Ponieważ wyniki trzech algorytmów na tej samej instancji są skorelowane
(trudność instancji wpływa na wszystkie trzy), zastosowano testy dla prób zależnych.
Istnienie jakichkolwiek różnic między algorytmami weryfikowano testem Friedmana
(rangowanie algorytmów w~obrębie każdej instancji, $k = 3$, korekta na remisy).
Różnice między parami algorytmów badano dwustronnym testem rangowanych znaków
Wilcoxona z~rozkładem dokładnym (brak remisów, $n = 20$), a~uzyskane
$p$-wartości korygowano metodą Holma ze względu na trzy porównania. Wielkość efektu
wyrażono deltą Cliffa $\delta \in [-1, 1]$ z~progami interpretacji
$|\delta| < 0{,}147$ (pomijalny), $< 0{,}33$ (mały), $< 0{,}474$ (średni), większe
(duży) oraz liczbą instancji, w~których dany algorytm uzyskał wyższą hiperobjętość.
Przyjęto poziom istotności $\alpha = 0{,}05$. Wszystkie obliczenia wykonano
skryptem \texttt{analyze\_comparison.py}, który nie korzysta z~zewnętrznych
bibliotek statystycznych --- testy zaimplementowano wprost, co ułatwia ich
audyt.

\section{Porównanie hiperobjętością}
\label{sec:analiza_hv}

Tabela~\ref{tab:hv_summary} zestawia statystyki znormalizowanej hiperobjętości
frontów końcowych w~eksperymencie głównym, a~rysunek~\ref{fig:hv_boxplot}
przedstawia jej rozkład po instancjach. Wartości dla poszczególnych instancji
zawiera tabela~\ref{tab:per_run} w~rozdziale~\ref{ch:wyniki}.

\begin{table}[htbp]
    \centering
    \caption{Znormalizowana hiperobjętość $\mathrm{HV}/\mathrm{HV}_{\max}$ frontów końcowych (20 instancji, $N = 100$, 500 iteracji)}
    \label{tab:hv_summary}
    \begin{tabular}{l c c c c c c}
        \hline
        Algorytm & Średnia $\pm$ odch. std. & Mediana & Min & Max & Średnia ranga & Zwycięstwa \\
        \hline
        NSGA-II & 0{,}728 $\pm$ 0{,}256 & 0{,}765 & 0{,}118 & 0{,}995 & 1{,}45 & 12 / 20 \\
        NSGA-III & 0{,}575 $\pm$ 0{,}270 & 0{,}546 & 0{,}088 & 0{,}977 & 1{,}65 & 8 / 20 \\
        MOGWO & 0{,}192 $\pm$ 0{,}154 & 0{,}129 & 0{,}045 & 0{,}641 & 2{,}90 & 0 / 20 \\
        \hline
    \end{tabular}
\end{table}

\begin{figure}[htbp]
    \centering
    \includegraphics[width=0.8\textwidth]{img/fig_hv_boxplot.pdf}
    \caption{Rozkład znormalizowanej hiperobjętości frontów końcowych po 20 instancjach (pudełko: kwartyle i~mediana; punkty: poszczególne instancje)}
    \label{fig:hv_boxplot}
\end{figure}

Test Friedmana odrzuca hipotezę o~braku różnic między algorytmami
($\chi^2_F = 24{,}70$, $p = 4{,}3 \cdot 10^{-6}$). Uporządkowanie algorytmów
według średniej rangi jest jednoznaczne: NSGA-II ($1{,}45$), NSGA-III ($1{,}65$),
MOGWO ($2{,}90$). Dwa algorytmy ewolucyjne dzielą między siebie wszystkie
zwycięstwa (odpowiednio $12$ i~$8$ instancji), natomiast MOGWO nie uzyskał
najwyższej hiperobjętości w~żadnej z~20 instancji, a~w~18 z~nich zajął ostatnie
miejsce.

Uwagę zwraca bardzo duży rozrzut wyników wszystkich algorytmów między instancjami
(odchylenie standardowe $0{,}15$--$0{,}27$ przy zakresie miary $[0, 1]$). Dla
NSGA-II i~NSGA-III rozkład jest wyraźnie asymetryczny: mediana NSGA-II ($0{,}765$)
przewyższa średnią, a~górny kwartyl sięga $0{,}955$, co oznacza, że w~jednej
czwartej instancji front NSGA-II niemal całkowicie dominuje fronty pozostałych
algorytmów. Jednocześnie w~instancjach 4 i~10 NSGA-II osiągnął zaledwie $0{,}19$
i~$0{,}12$, podczas gdy NSGA-III uzyskał tam $0{,}64$ i~$0{,}95$; odwrotną sytuację
obserwuje się w~instancjach 16 i~19 (NSGA-III $0{,}30$ i~$0{,}09$ wobec $0{,}93$
i~$0{,}97$ dla NSGA-II). Przewaga jednego z~algorytmów ewolucyjnych nad drugim
nie jest więc cechą systematyczną, lecz zależy od instancji --- obserwacja ta
znajduje potwierdzenie w~testach dla par.


\section{Testy dla par algorytmów}
\label{sec:analiza_pary}

Wyniki testów Wilcoxona dla trzech par algorytmów zestawiono
w~tabeli~\ref{tab:pairwise}. Statystyka $W$ oznacza mniejszą z~sum rang różnic
dodatnich i~ujemnych; mediana $\Delta$ to mediana różnicy
$\mathrm{HV}/\mathrm{HV}_{\max}$ algorytmu $A$ i~$B$ w~tej samej instancji.

\begin{table}[htbp]
    \centering
    \caption{Testy Wilcoxona dla par (poprawka Holma) i~wielkość efektu $\delta$ Cliffa dla $\mathrm{HV}/\mathrm{HV}_{\max}$}
    \label{tab:pairwise}
    \begin{tabular}{l c c c c c l c}
        \hline
        Para $A$ vs $B$ & Mediana $\Delta$ & $W$ & $p$ & $p_{\text{Holm}}$ & $\delta$ Cliffa & Efekt & $A$ lepszy / $B$ lepszy \\
        \hline
        NSGA-II vs NSGA-III & 0{,}261 & 62{,}0 & $0{,}114$ & $0{,}114$ & 0{,}32 & mały & 12 / 8 \\
        NSGA-II vs MOGWO & 0{,}603 & 4{,}0 & $< 0{,}001$ & $< 0{,}001$ & 0{,}88 & duży & 19 / 1 \\
        NSGA-III vs MOGWO & 0{,}346 & 4{,}0 & $< 0{,}001$ & $< 0{,}001$ & 0{,}81 & duży & 19 / 1 \\
        \hline
    \end{tabular}
\end{table}

Obraz jest dwudzielny. Oba algorytmy ewolucyjne są istotnie lepsze od MOGWO
($p_{\text{Holm}} = 4 \cdot 10^{-5}$ w~obu porównaniach), a~efekt jest duży
($\delta = 0{,}88$ i~$0{,}81$): MOGWO przegrywa z~każdym z~nich w~19 z~20 instancji,
a~mediana różnicy hiperobjętości wynosi $0{,}60$ względem NSGA-II i~$0{,}35$
względem NSGA-III. Jedyną instancją, w~której MOGWO pokonał NSGA-II, jest
instancja 10 ($0{,}512$ wobec $0{,}118$); NSGA-III pokonał tam jednak oba
algorytmy ($0{,}949$).

Różnica między NSGA-II a~NSGA-III nie osiąga istotności statystycznej
($p = 0{,}114$; $W = 62$ przy sumach rang $W^{+} = 148$ i~$W^{-} = 62$), a~efekt
jest mały ($\delta = 0{,}32$). NSGA-II uzyskał wyższą hiperobjętość w~12
instancjach, NSGA-III w~8; mediana różnicy ($0{,}26$) jest wprawdzie wyraźnie
dodatnia, lecz rozkład różnic jest szeroki i~obejmuje wartości od $-0{,}83$
(instancja 10) do $+0{,}88$ (instancja 19). Wynik należy interpretować jako brak
dowodu na systematyczną przewagę któregokolwiek z~algorytmów ewolucyjnych przy
$n = 20$, a~nie jako dowód ich równoważności: przy zaobserwowanej wielkości efektu
wykazanie różnicy na poziomie $\alpha = 0{,}05$ wymagałoby znacząco większej
liczby instancji.

\section{Zachowanie w~przestrzeni kryteriów}
\label{sec:analiza_kryteria}

Hiperobjętość agreguje trzy kryteria w~jedną liczbę; tabela~\ref{tab:objectives}
i~rysunek~\ref{fig:best_objectives} pokazują, z~czego ta różnica wynika.
Dla każdej instancji wyznaczono najlepszą wartość każdego kryterium na froncie
końcowym (uśrednioną następnie po instancjach), udział instancji, w~których
wszystkie drony dotarły do celu w~horyzoncie symulacji, udział rozwiązań
o~zerowym ryzyku radarowym na froncie oraz rozmiar frontu.

\begin{table}[htbp]
    \centering
    \caption{Średnie (po instancjach) najlepsze wartości kryteriów na froncie, udział instancji z~wszystkimi dronami u~celu, udział rozwiązań o~zerowym ryzyku, rozmiar frontu i~czas obliczeń}
    \label{tab:objectives}
    \begin{tabular}{l c c c c c c}
        \hline
        Algorytm & Min. makespan & Min. energia & Wszystkie drony u celu & Udział ryzyka $= 0$ & $|$front$|$ & Czas [s] \\
        \hline
        NSGA-II & 455{,}4 & 804{,}0 & 100\% & 31{,}1\% & 26{,}8 & 3{,}60 \\
        NSGA-III & 528{,}0 & 857{,}8 & 95\% & 27{,}3\% & 27{,}1 & 3{,}65 \\
        MOGWO & 797{,}7 & 931{,}7 & 70\% & 25{,}9\% & 56{,}0 & 6{,}38 \\
        \hline
    \end{tabular}
\end{table}

\begin{figure}[htbp]
    \centering
    \includegraphics[width=\textwidth]{img/fig_best_objectives.pdf}
    \caption{Najlepszy makespan i~najlepsza energia na froncie końcowym oraz udział rozwiązań o~zerowym ryzyku radarowym --- rozkłady po 20 instancjach}
    \label{fig:best_objectives}
\end{figure}

Kryterium, które najsilniej różnicuje algorytmy, jest makespan. Horyzont
symulacji wynosi $900$ tików, a~drony, które nie dotarły do celu, otrzymują
makespan większy od horyzontu. NSGA-II we wszystkich 20 instancjach doprowadził
wszystkie drony do celu, osiągając najlepszy makespan średnio $455$ tików
(mediana $408$, minimum $288$), NSGA-III w~19 instancjach (średnio $528$),
natomiast MOGWO tylko w~14 z~20 instancji, a~jego najlepszy makespan wynosił
średnio $798$ tików (mediana $852$) --- w~typowej instancji wataha nie zdołała
znaleźć trasy, która kończy misję wyraźnie przed horyzontem. Analogiczne, choć
słabsze, uporządkowanie widać w~energii: najlepsze rozwiązania MOGWO zużywają
średnio o~$16\%$ więcej energii niż NSGA-II.

Kryterium ryzyka radarowego nie różnicuje algorytmów pod względem wartości
minimalnej: każdy z~nich w~każdej instancji znalazł trasy całkowicie omijające
strefy radarów (ryzyko $0$). Różnice dotyczą udziału takich rozwiązań we froncie
($31\%$ dla NSGA-II, $27\%$ dla NSGA-III, $26\%$ dla MOGWO) oraz --- co
istotniejsze --- jakości kompromisów: trasa bezpieczna jest użyteczna wtedy, gdy
jednocześnie kończy misję w~rozsądnym czasie, a~to właśnie na osi makespanu MOGWO
pozostaje w~tyle. Warto odnotować dwumodalny charakter tego wskaźnika dla MOGWO
(odchylenie standardowe $25$ punktów procentowych wobec $12$ dla algorytmów
ewolucyjnych): w~kilku instancjach archiwum MOGWO składało się niemal wyłącznie
z~tras bezpiecznych, lecz powolnych.

Rozmiar frontu obnaża pozorną zaletę MOGWO: jego archiwum liczy średnio $56$
rozwiązań (do $100$, czyli pełna pojemność), dwukrotnie więcej niż fronty
algorytmów ewolucyjnych, lecz przy ponad trzykrotnie mniejszej hiperobjętości.
Liczba rozwiązań niezdominowanych nie jest więc miarą jakości frontu --- archiwum
MOGWO gromadzi wiele wzajemnie niezdominowanych, ale globalnie słabych tras.
Rysunek~\ref{fig:fronts_projections} ilustruje to na przykładzie instancji 1:
front MOGWO (hiperobjętość $0{,}08$) skupia się w~obszarze dużego makespanu,
a~fronty NSGA-II ($0{,}58$) i~NSGA-III ($0{,}86$) rozciągają się w~kierunku
krótszych i~oszczędniejszych tras, przy czym NSGA-III pokrywa szerszy zakres
kompromisu czas--energia.

\begin{figure}[htbp]
    \centering
    \includegraphics[width=\textwidth]{img/fig_fronts_projections.pdf}
    \caption{Fronty końcowe trzech algorytmów w~instancji 1 --- rzuty na pary kryteriów}
    \label{fig:fronts_projections}
\end{figure}

\section{Zbieżność}
\label{sec:analiza_zbieznosc}

Rysunek~\ref{fig:convergence} przedstawia średnią (po instancjach) znormalizowaną
hiperobjętość bieżącego zbioru niezdominowanego w~funkcji numeru iteracji,
a~tabela~\ref{tab:convergence} --- jej wartości po kolejnych ułamkach budżetu
oraz medianę liczby iteracji potrzebnych do osiągnięcia $90\%$ hiperobjętości
końcowej.

\begin{figure}[htbp]
    \centering
    \includegraphics[width=0.85\textwidth]{img/fig_convergence.pdf}
    \caption{Krzywe zbieżności: średnia $\mathrm{HV}/\mathrm{HV}_{\max}$ bieżącego zbioru niezdominowanego (pasmo: $\pm$ odchylenie standardowe po 20 instancjach)}
    \label{fig:convergence}
\end{figure}

\begin{table}[htbp]
    \centering
    \caption{Średnia $\mathrm{HV}/\mathrm{HV}_{\max}$ bieżącego zbioru niezdominowanego po ułamku budżetu (500 iteracji) oraz mediana liczby iteracji do osiągnięcia $90\%$ hiperobjętości końcowej}
    \label{tab:convergence}
    \begin{tabular}{l c c c c c c}
        \hline
        Algorytm & 20\% & 40\% & 60\% & 80\% & 100\% & Iteracje do 90\% HV$_{\text{końc.}}$ \\
        \hline
        NSGA-II & 0{,}121 & 0{,}304 & 0{,}458 & 0{,}619 & 0{,}728 & 425 \\
        NSGA-III & 0{,}078 & 0{,}218 & 0{,}354 & 0{,}469 & 0{,}575 & 430 \\
        MOGWO & 0{,}023 & 0{,}065 & 0{,}152 & 0{,}178 & 0{,}192 & 385 \\
        \hline
    \end{tabular}
\end{table}

Krzywe algorytmów ewolucyjnych są niemal liniowe i~do końca budżetu nie
wykazują nasycenia: w~ostatnich $20\%$ iteracji NSGA-II zyskał $0{,}11$,
a~NSGA-III $0{,}11$ hiperobjętości, a~mediana liczby iteracji do $90\%$ wartości
końcowej ($425$ i~$430$) leży blisko końca budżetu. Oznacza to, że przy $500$
iteracjach oba algorytmy są wciąż w~fazie poprawy i~dłuższe obliczenia
przyniosłyby dalszy wzrost jakości --- co potwierdza analiza budżetu
w~podrozdziale~\ref{sec:analiza_budzet}. NSGA-II przez cały przebieg utrzymuje
przewagę nad NSGA-III o~stałej w~przybliżeniu proporcji (ok. $0{,}8$ wartości
NSGA-II), bez zmiany uporządkowania w~żadnym punkcie.

Krzywa MOGWO ma inny kształt. Po powolnym starcie (hiperobjętość $0{,}02$ po
$20\%$ budżetu wobec $0{,}12$ dla NSGA-II) następuje faza najszybszego wzrostu
między $40\%$ a~$60\%$ budżetu, po czym krzywa wyraźnie się wypłaszcza: w~ostatnich
$40\%$ iteracji przyrost wynosi zaledwie $0{,}04$ (z~$0{,}152$ do~$0{,}192$),
a~$90\%$ wartości końcowej MOGWO osiąga najwcześniej ze wszystkich algorytmów
(mediana $385$ iteracji). Wczesne osiągnięcie plateau nie jest tu zaletą, lecz
objawem stagnacji: algorytm przestaje znajdować nowe rozwiązania niezdominowane,
zanim wyczerpie budżet. Przyczyny tego zjawiska omówiono
w~podrozdziale~\ref{sec:analiza_interpretacja}.

\section{Wrażliwość na budżet obliczeniowy i~czas obliczeń}
\label{sec:analiza_budzet}

Ponieważ normalizacja hiperobjętości per instancja uniemożliwia porównanie
wartości uzyskanych przy różnych budżetach (każdy budżet wyznaczałby własny
punkt nadir), w~tej części analizy wartości kryteriów znormalizowano do granic
wspólnych dla instancji po \emph{wszystkich} budżetach i~algorytmach. Dzięki temu
wzrost hiperobjętości wraz z~budżetem odzwierciedla rzeczywistą poprawę frontów.
Wyniki dla $10$ instancji i~czterech budżetów przedstawia tabela~\ref{tab:budget}
oraz rysunek~\ref{fig:budget}.

\begin{table}[htbp]
    \centering
    \caption{Wrażliwość na budżet obliczeniowy: średnia $\pm$ odch. std. $\mathrm{HV}/\mathrm{HV}_{\max}$ przy normalizacji wspólnej dla instancji po wszystkich budżetach (10 instancji na budżet); zwycięstwa w~kolejności NSGA-II / NSGA-III / MOGWO}
    \label{tab:budget}
    \begin{tabular}{l c c c c c}
        \hline
        Iteracje & NSGA-II & NSGA-III & MOGWO & Zwycięstwa & $p$ Friedmana \\
        \hline
        100 & 0{,}152 $\pm$ 0{,}104 & 0{,}169 $\pm$ 0{,}089 & 0{,}067 $\pm$ 0{,}015 & 3 / 7 / 0 & $0{,}001$ \\
        250 & 0{,}390 $\pm$ 0{,}195 & 0{,}385 $\pm$ 0{,}200 & 0{,}168 $\pm$ 0{,}141 & 5 / 4 / 1 & $0{,}061$ \\
        500 & 0{,}586 $\pm$ 0{,}255 & 0{,}569 $\pm$ 0{,}183 & 0{,}243 $\pm$ 0{,}184 & 5 / 5 / 0 & $0{,}002$ \\
        1000 & 0{,}806 $\pm$ 0{,}242 & 0{,}803 $\pm$ 0{,}150 & 0{,}442 $\pm$ 0{,}353 & 5 / 3 / 2 & $0{,}150$ \\
        \hline
    \end{tabular}
\end{table}

\begin{figure}[htbp]
    \centering
    \includegraphics[width=0.8\textwidth]{img/fig_budget.pdf}
    \caption{Znormalizowana hiperobjętość (wspólna normalizacja po budżetach) w~funkcji liczby iteracji --- średnia $\pm$ odchylenie standardowe po 10 instancjach}
    \label{fig:budget}
\end{figure}

Trzy obserwacje wynikają z~tych danych. Po pierwsze, jakość frontów wszystkich
algorytmów rośnie monotonicznie i~bez nasycenia w~całym badanym zakresie:
dziesięciokrotne zwiększenie budżetu podnosi hiperobjętość NSGA-II z~$0{,}15$ do
$0{,}81$, a~MOGWO z~$0{,}07$ do $0{,}44$. Po drugie, uporządkowanie algorytmów
ewolucyjnych zależy od budżetu: przy $100$ iteracjach NSGA-III jest najlepszy
w~7 z~10 instancji (średnia ranga $1{,}30$, test Friedmana $p = 0{,}001$), przy
$250$--$500$ iteracjach oba algorytmy są praktycznie równoważne, a~przy $1000$
iteracjach ich średnie hiperobjętości różnią się o~$0{,}003$. Mechanizm niszowania
z~punktami referencyjnymi daje zatem przewagę w~początkowej fazie --- równomierne
rozpychanie frontu jest skuteczne, gdy populacja jest jeszcze daleko od frontu ---
lecz przewaga ta zanika wraz z~budżetem. Po trzecie, MOGWO pozostaje najsłabszy
przy każdym budżecie, ale jego wynik przy $1000$ iteracjach ($0{,}44 \pm 0{,}35$)
kryje największy rozrzut ze wszystkich pomiarów: w~instancjach 5 i~10 archiwum
MOGWO osiągnęło hiperobjętość $0{,}999$ i~$0{,}998$, dominując niemal cały obszar
rozpięty przez pozostałe fronty (NSGA-II uzyskał tam $0{,}28$ i~$0{,}13$), podczas
gdy w~pozostałych ośmiu instancjach nie przekroczył $0{,}33$. Dyskretny wariant
GWO jest więc zdolny do znalezienia bardzo dobrych tras, lecz czyni to w~sposób
niepowtarzalny --- co dla zastosowań praktycznych jest wadą równie poważną, jak
niska średnia.

Czas obliczeń (tabela~\ref{tab:objectives}, ostatnia kolumna) rośnie liniowo
z~budżetem i~dla algorytmów ewolucyjnych jest niemal identyczny ($3{,}6$ s
i~$3{,}7$ s na $50\,000$ ocen, tj. ok.~$72\,\mu$s na ocenę trasy), co potwierdza,
że koszt zdominowany jest przez symulację, a~nie przez operatory selekcji. MOGWO
potrzebuje $1{,}8$ raza więcej czasu ($6{,}4$ s przy tym samym budżecie ocen), co
wynika z~dyskretnej aktualizacji pozycji --- dla każdego z~$900$ genów każdego
wilka losowane są trzy współczynniki $A_k$ i~trzy propozycje --- oraz z~utrzymania
archiwum (obliczanie odległości zatłoczenia przy każdej redukcji). Ponieważ
porównanie prowadzono przy równej liczbie ocen, różnica ta działa na niekorzyść
MOGWO dodatkowo, gdyby za miarę budżetu przyjąć czas procesora.

\section{Interpretacja różnic}
\label{sec:analiza_interpretacja}

\subsection{Dlaczego MOGWO ustępuje algorytmom ewolucyjnym}
\label{sec:analiza_mogwo}
Zaobserwowana różnica jest duża, systematyczna i~obecna przy każdym budżecie,
należy więc szukać jej przyczyn w~samej konstrukcji dyskretnego operatora
aktualizacji pozycji (podrozdział~\ref{sec:mogwo_discrete}), a~nie w~losowości.
Na podstawie zebranych danych można wskazać cztery współdziałające mechanizmy.

\begin{enumerate}
    \item \textbf{Epistaza kodowania a~operator genowy.} Znaczenie genu $g_{d,s}$
    zależy od pozycji, w~której dron znajduje się w~kroku $s$, a~ta wynika
    z~całego prefiksu trasy. Głosowanie genowe przenosi od liderów pojedyncze
    kody ruchów do ustalonych pozycji w~chromosomie, ignorując kontekst: kod
    ,,w~górę i~w~prawo'', który u~lidera omijał szczyt, u~wilka o~innym prefiksie
    prowadzi w~przypadkową stronę. Krzyżowanie jednopunktowe stosowane w~NSGA
    wymienia natomiast spójne sufiksy tras, zachowując ich wewnętrzną strukturę.
    Operator MOGWO działa w~praktyce jak krzyżowanie jednorodne z~trzema rodzicami
    połączone z~silną mutacją --- operator o~wysokiej destrukcyjności dla kodowań
    z~zależnościami pozycyjnymi.
    \item \textbf{Brak elitaryzmu w~watasze.} W~każdej iteracji wszystkie $N$
    wilków są zastępowane nowymi pozycjami; dobre rozwiązania przeżywają
    wyłącznie w~archiwum i~oddziałują na populację tylko przez trzech liderów.
    W~NSGA-II/III cała elita ($N$ najlepszych z~$2N$) przechodzi do następnego
    pokolenia i~uczestniczy w~reprodukcji, co daje znacznie silniejszą presję
    selekcyjną. Efekt widać w~makespanie: wataha w~$30\%$ instancji nie
    utrwaliła nawet tras doprowadzających wszystkie drony do celu.
    \item \textbf{Wygaszanie eksploracji przez harmonogram parametru $a$.} Gdy
    $a \rightarrow 0$, warunek $|A_k| < 1$ jest spełniony niemal zawsze, więc
    każdy gen jest przejmowany od liderów, a~prawdopodobieństwo mutacji
    $0{,}05\,a$ spada do zera. Wataha zbiega do kombinacji trzech liderów,
    różnorodność zanika i~archiwum przestaje otrzymywać nowe rozwiązania ---
    dokładnie to pokazuje wypłaszczenie krzywej zbieżności po $60\%$ budżetu
    (podrozdział~\ref{sec:analiza_zbieznosc}). W~ciągłym GWO ten sam harmonogram
    jest pożądany, bo końcowa faza służy lokalnemu doprecyzowaniu rozwiązania;
    w~przestrzeni dyskretnej nie istnieje ,,małe przesunięcie'' w~kierunku lidera
    i~faza ta degeneruje się do kopiowania.
    \item \textbf{Wybór liderów niezależny od jakości.} Ruletka po odległości
    zatłoczenia preferuje rzadko zaludnione obszary archiwum bez względu na to,
    czy reprezentują one wartościowe kompromisy. Gdy archiwum zawiera skrajne
    rozwiązania o~minimalnej energii, lecz z~dronami, które nie dotarły do celu,
    takie rozwiązania są wybierane na liderów szczególnie często --- i~ciągną
    watahę w~stronę tras bezpiecznych, ale bezużytecznych. Tłumaczy to dwumodalny
    rozkład udziału rozwiązań o~zerowym ryzyku w~archiwach MOGWO.
\end{enumerate}

Mechanizmy te nie są nieusuwalne. Wyniki z~instancji 5 i~10 przy $1000$
iteracjach dowodzą, że gdy wataha trafi na właściwy region przestrzeni, archiwum
potrafi wypełnić go lepiej niż NSGA. Naturalnymi kierunkami poprawy są: wymiana
od liderów spójnych segmentów tras zamiast pojedynczych genów, zachowanie części
archiwum w~watasze (elitaryzm) oraz dolne ograniczenie prawdopodobieństwa mutacji
niezależne od $a$. Weryfikacja tych modyfikacji wykracza jednak poza zakres
niniejszej pracy, której celem było porównanie algorytmów w~postaciach opisanych
w~rozdziale~\ref{ch:sota}.

\subsection{NSGA-II a~NSGA-III}
\label{sec:analiza_nsga}
Brak istotnej różnicy między NSGA-II a~NSGA-III jest zgodny z~wynikami
Ishibuchiego i~in.~\cite{ishibuchi2016nsga}, którzy dla dwóch i~trzech kryteriów
nie stwierdzili systematycznej przewagi NSGA-III. Rozważany problem ma trzy
kryteria, a~więc leży poza zakresem, dla którego zaprojektowano punkty
referencyjne: sortowanie niezdominowane nadal skutecznie różnicuje populację
(fronty końcowe liczą średnio $27$ rozwiązań przy populacji $100$, więc presja
selekcyjna rang jest wysoka), a~odległość zatłoczenia wystarcza do utrzymania
rozpiętości frontu. Przewaga NSGA-III przy małym budżecie ($100$ iteracji)
wskazuje, że niszowanie z~punktami referencyjnymi szybciej rozprowadza populację
wzdłuż frontu, natomiast turniej binarny z~operatorem porównania zatłoczonego
w~NSGA-II zapewnia silniejszą presję selekcyjną w~dalszych fazach, gdy liczy się
głównie zbieżność. Oba mechanizmy prowadzą przy $1000$ iteracjach do frontów
o~praktycznie tej samej jakości. Dla praktyki oznacza to, że w~zadaniu tej klasy
wybór między NSGA-II a~NSGA-III nie jest krytyczny, a~decydujące znaczenie ma
budżet obliczeniowy.

\section{Ocena przydatności w~warunkach niestacjonarnych}
\label{sec:analiza_niestacjonarnosc}

Niestacjonarność środowiska jest w~symulatorze realizowana przez zmienne w~czasie
pole wiatru (podrozdział~\ref{sec:niestacjonarnosc}): koszt energetyczny tego
samego kroku zależy od tiku, w~którym zostaje wykonany, więc przesunięcie trasy
w~czasie --- np. wskutek odrzucenia wcześniejszego ruchu --- zmienia jej koszt.
Wszystkie trzy algorytmy pracują w~trybie planowania z~pełną znajomością modelu
wiatru, zatem niestacjonarność nie stawia przed nimi wymogu adaptacji w~locie,
lecz utrudnia zadanie przez uzależnienie wartości genów od ich położenia w~czasie.
Jest to ten sam rodzaj zależności pozycyjnej, który w~podrozdziale
\ref{sec:analiza_mogwo} wskazano jako główną słabość operatora genowego MOGWO:
algorytm, którego operator ignoruje kontekst genu, jest szczególnie wrażliwy na
środowisko, w~którym kontekst ten obejmuje również czas. Operatory NSGA, wymieniające
spójne fragmenty tras, zachowują zarówno kontekst przestrzenny, jak i~czasowy.

Przeprowadzone eksperymenty nie pozwalają jednak ilościowo wyodrębnić wpływu
samej niestacjonarności, ponieważ wszystkie instancje zawierały zmienny wiatr.
Rozstrzygnięcie, o~ile wiatr zmienny w~czasie pogarsza wyniki względem wiatru
stałego, wymagałoby eksperymentu ablacyjnego z~wyłączonymi podmuchami przy
niezmienionych pozostałych parametrach; jest to najprostsze i~najbardziej
wartościowe rozszerzenie niniejszej analizy.

\section{Ograniczenia analizy}
\label{sec:analiza_ograniczenia}

\begin{itemize}
    \item \textbf{Względny charakter metryki.} Normalizacja per instancja sprawia,
    że $\mathrm{HV}/\mathrm{HV}_{\max}$ mierzy jakość frontu względem pozostałych
    algorytmów w~tej samej instancji, a~nie względem nieznanego frontu
    optymalnego. Wartość $0{,}99$ oznacza dominację nad konkurentami, a~nie
    bliskość optimum.
    \item \textbf{Wybór punktu referencyjnego.} Margines $10\%$ jest konwencją;
    większy margines zwiększa wagę skrajnych rozwiązań frontu. Wrażliwości
    wyników na ten parametr nie badano.
    \item \textbf{Liczebność próby.} Dwadzieścia instancji wystarczyło do
    wykazania dużego efektu (MOGWO), ale nie małego (NSGA-II vs NSGA-III);
    w~serii budżetowej $n = 10$ dodatkowo ogranicza moc testów.
    \item \textbf{Jedna rodzina scenariuszy.} Wszystkie instancje mają ten sam
    rozmiar mapy, liczbę dronów i~radarów oraz ten sam wiatr bazowy; różnią się
    jedynie rzeźbą terenu. Uogólnienie na inne konfiguracje wymaga osobnych
    eksperymentów.
    \item \textbf{Specyfika adaptacji MOGWO.} Wnioski dotyczą opisanego
    w~pracy dyskretnego wariantu z~głosowaniem genowym i~archiwum opartym na
    odległości zatłoczenia, a~nie algorytmu MOGWO jako takiego; inna dyskretyzacja
    operatora może dać odmienne wyniki.
    \item \textbf{Pomiar czasu.} Czasy obliczeń mierzono na współdzielonej
    maszynie; ich wartości bezwzględne są orientacyjne, miarodajne są proporcje.
\end{itemize}
